% House preamble for every reward-lens TikZ figure.
% The figure file sets \documentclass{standalone} and, before it is \input, the build
% sets \rltheme (light|dark). This file loads the matching palette and defines the shared
% visual vocabulary, so consistency across figures comes from here, not from discipline.
%
% Engine: lualatex or xelatex gives the real Inter (matches the site). pdflatex falls back
% to Fira Sans, which is close. Either way the figures do not speak Computer Modern.

\usepackage{iftex}
\ifLuaTeX
  \usepackage{fontspec}\setsansfont{Inter}[BoldFont={Inter SemiBold}, ItalicFont={Inter Italic}]
\else\ifXeTeX
  \usepackage{fontspec}\setsansfont{Inter}[BoldFont={Inter SemiBold}, ItalicFont={Inter Italic}]
\else
  \usepackage[T1]{fontenc}\usepackage[sfdefault]{FiraSans}
\fi\fi
\renewcommand{\familydefault}{\sfdefault}

\usepackage{amsmath}
\usepackage{tikz}
\usetikzlibrary{
  positioning, calc, arrows.meta, fit, backgrounds, shapes.geometric,
  shapes.misc, decorations.pathreplacing, decorations.pathmorphing,
  shadows.blur, patterns, math
}

% AUTO-GENERATED from palette.json by emit_palette.py. Do not edit by hand.
% Each figure sets \rltheme to light or dark before \input; the preamble loads this file.
\usepackage{etoolbox}
\providecommand{\rltheme}{light}%
\ifdefstring{\rltheme}{dark}{%
  \definecolor{rlInk}{HTML}{E6E8EB}%
  \definecolor{rlInkTwo}{HTML}{B7BDC7}%
  \definecolor{rlMuted}{HTML}{9AA0AA}%
  \definecolor{rlLine}{HTML}{3A3D44}%
  \definecolor{rlSurface}{HTML}{1D1E22}%
  \definecolor{rlSurfaceTwo}{HTML}{26272C}%
  \definecolor{rlBaseline}{HTML}{55565E}%
  \definecolor{rlTierObs}{HTML}{2DD4BF}%
  \definecolor{rlTierCausal}{HTML}{F59E0B}%
  \definecolor{rlTierVuln}{HTML}{FB7185}%
  \definecolor{rlTrustExp}{HTML}{64748B}%
  \definecolor{rlTrustCal}{HTML}{60A5FA}%
  \definecolor{rlTrustReg}{HTML}{A78BFA}%
  \definecolor{rlTrustAdj}{HTML}{22C55E}%
  \definecolor{rlSignal}{HTML}{E879F9}%
  \definecolor{rlGauge}{HTML}{22D3EE}%
  \definecolor{rlChosen}{HTML}{5598E7}%
  \definecolor{rlRejected}{HTML}{F0716F}%
  \definecolor{rlSeriesB}{HTML}{F98A5C}%
  \definecolor{rlAccent}{HTML}{FB8C5A}%
  \definecolor{rlSeqLo}{HTML}{0D366B}%
  \definecolor{rlSeqHi}{HTML}{CDE2FB}%
  \definecolor{rlDivMid}{HTML}{2A2B31}%
}{%
  \definecolor{rlInk}{HTML}{16181D}%
  \definecolor{rlInkTwo}{HTML}{4B5563}%
  \definecolor{rlMuted}{HTML}{6B7280}%
  \definecolor{rlLine}{HTML}{D7DAE0}%
  \definecolor{rlSurface}{HTML}{FFFFFF}%
  \definecolor{rlSurfaceTwo}{HTML}{F4F5F7}%
  \definecolor{rlBaseline}{HTML}{C3C2B7}%
  \definecolor{rlTierObs}{HTML}{0D9488}%
  \definecolor{rlTierCausal}{HTML}{B45309}%
  \definecolor{rlTierVuln}{HTML}{BE123C}%
  \definecolor{rlTrustExp}{HTML}{94A3B8}%
  \definecolor{rlTrustCal}{HTML}{3B82F6}%
  \definecolor{rlTrustReg}{HTML}{7C3AED}%
  \definecolor{rlTrustAdj}{HTML}{15803D}%
  \definecolor{rlSignal}{HTML}{A21CAF}%
  \definecolor{rlGauge}{HTML}{0891B2}%
  \definecolor{rlChosen}{HTML}{2A78D6}%
  \definecolor{rlRejected}{HTML}{E34948}%
  \definecolor{rlSeriesB}{HTML}{EB6834}%
  \definecolor{rlAccent}{HTML}{E85D2C}%
  \definecolor{rlSeqLo}{HTML}{CDE2FB}%
  \definecolor{rlSeqHi}{HTML}{0D366B}%
  \definecolor{rlDivMid}{HTML}{F0EFEC}%
}%
 % defines the rl* colors for the current \rltheme

% digit-name aliases, so both rlInkTwo and rlInk2 resolve
\colorlet{rlInk2}{rlInkTwo}
\colorlet{rlSurface2}{rlSurfaceTwo}

% ---- global defaults -------------------------------------------------------
\tikzset{
  >={Stealth[length=2.4mm, width=1.8mm]},
  font=\sffamily,
  every node/.append style={text=rlInk},
}

% ---- chips: the three-chip vocabulary (tier / gauge status / works-on) ------
% Outline chip reads correctly on both themes because the role color is a mid-tone.
\tikzset{
  rlchip/.style 2 args={
    draw=#1, text=#1, line width=0.5pt, rounded corners=2.5pt,
    inner xsep=5pt, inner ysep=2.6pt, font=\sffamily\bfseries\scriptsize,
    label/.expanded=#2,
  },
  rlchip/.default={rlMuted}{},
  % filled variant, surface-colored text so it contrasts in both themes
  rlchipfill/.style={
    draw=#1, fill=#1, text=rlSurface, rounded corners=2.5pt,
    inner xsep=5pt, inner ysep=2.6pt, font=\sffamily\bfseries\scriptsize,
  },
}
% simple outline chip taking one color (most common)
\newcommand{\rlchipnode}[3][rlMuted]{\node[draw=#1, text=#1, line width=0.6pt,
  rounded corners=2.5pt, inner xsep=5pt, inner ysep=2.6pt,
  font=\sffamily\bfseries\scriptsize] (#2) {#3};}

% ---- cards and panels ------------------------------------------------------
\tikzset{
  rlcard/.style={
    draw=rlLine, fill=rlSurface, line width=0.7pt, rounded corners=4pt,
    blur shadow={shadow blur steps=7, shadow xshift=0pt, shadow yshift=-1pt,
      shadow opacity=18, shadow blur radius=3pt},
  },
  rlpanel/.style={draw=rlLine, fill=rlSurfaceTwo, line width=0.6pt, rounded corners=3pt},
  rlfield/.style={font=\sffamily\scriptsize, text=rlMuted, align=left},
  rlvalue/.style={font=\sffamily\bfseries, text=rlInk},
  rllabel/.style={font=\sffamily\footnotesize, text=rlInk2},
  rlnote/.style={font=\sffamily\scriptsize\itshape, text=rlMuted},
}

% ---- lines, axes, vectors --------------------------------------------------
\tikzset{
  rlaxis/.style={draw=rlLine, line width=0.7pt, -{Stealth[length=1.8mm]}},
  rlbaseline/.style={draw=rlBaseline, line width=0.6pt, dash pattern=on 2pt off 2pt},
  rlflow/.style={draw=rlInk2, line width=0.9pt, -{Stealth[length=2.4mm]}},
  rlvec/.style={line width=1.3pt, -{Stealth[length=2.8mm, width=2.2mm]}},
  rldrop/.style={draw=rlMuted, line width=0.5pt, dash pattern=on 1.6pt off 1.6pt},
}

% ---- gate glyph: a narrow keyhole the evidence must pass to climb a rung ----
% \rlgate{name}{fill color}{label}  draws a small rounded gate node.
\newcommand{\rlgate}[3]{%
  \node[draw=#2, fill=rlSurface, line width=0.9pt, rounded corners=2pt,
        minimum width=6.5mm, minimum height=6.5mm, inner sep=0pt] (#1) {};
  \node[text=#2, font=\sffamily\bfseries\footnotesize] at (#1) {#3};%
}

% ---- answer-key glyph (calibration): a tiny key ----------------------------
\newcommand{\rlkey}[2][rlTrustAdj]{%
  \begin{scope}[shift={(#2)}, rotate=45]
    \draw[fill=#1, draw=#1] (0,0) circle (1.6mm);
    \fill[rlSurface] (0,0) circle (0.6mm);
    \draw[#1, line width=1.1pt] (0,-1.6mm) -- (0,-4.4mm);
    \draw[#1, line width=1.1pt] (0,-3.0mm) -- (1.1mm,-3.0mm);
    \draw[#1, line width=1.1pt] (0,-4.0mm) -- (0.9mm,-4.0mm);
  \end{scope}%
}

% ---- frozen-study seal: a small notched stamp ------------------------------
\newcommand{\rlseal}[2][rlTrustReg]{%
  \node[star, star points=12, star point ratio=0.9, draw=#1, fill=#1,
        text=rlSurface, minimum size=9mm, inner sep=0pt,
        font=\sffamily\bfseries\tiny] at (#2) {\#hash};%
}
